\ifcsname LRACalculusCapstoneRendered\endcsname
\else
\expandafter\gdef\csname LRACalculusCapstoneRendered\endcsname{}

\newpage
\phantomsection
\label{cap:computational-calculus}

\begin{tcolorbox}[
  colback=gray!6,
  colframe=gray!40,
  arc=2pt,
  left=8pt, right=8pt, top=6pt, bottom=6pt,
  title={\small\textbf{Capstone Problem}},
  fonttitle=\small\bfseries
]
\textbf{Problem.}
TODO: state one synthesizing problem for Calculus, using only concepts at or
before this chapter in the registry.
\end{tcolorbox}

\begin{remark*}[Dependency ceiling]
The capstone may use only results routed at or before this chapter.
\end{remark*}

% =========================================================
% Migrated exercises from limits-and-continuity.tex
% =========================================================

% =========================================================
% Calculus --- Limits: Exercises
% =========================================================

\subsection*{Exercises: Limits}

% ---------------------------------------------------------
\subsubsection*{Direct Evaluation and Limit Laws}

\begin{enumerate}

\item Evaluate the following limits using direct substitution or
      the limit laws.
      \begin{enumerate}
        \item $\displaystyle\lim_{x \to 3}\ (2x^2 - 5x + 1)$
        \item $\displaystyle\lim_{x \to -2}\ \frac{x^2 + 3x + 2}{x + 1}$
        \item $\displaystyle\lim_{x \to 4}\ \sqrt{3x - 2}$
        \item $\displaystyle\lim_{x \to 0}\ \frac{\sin(x)}{x}$
              \quad (\textit{state the result; no proof required})
      \end{enumerate}

\item Evaluate each limit by factoring and cancelling.
      \begin{enumerate}
        \item $\displaystyle\lim_{x \to 2}\ \frac{x^2 - 4}{x - 2}$
        \item $\displaystyle\lim_{x \to -3}\ \frac{x^2 + 5x + 6}{x + 3}$
        \item $\displaystyle\lim_{x \to 1}\ \frac{x^3 - 1}{x - 1}$
      \end{enumerate}

\item Evaluate by rationalising.
      \begin{enumerate}
        \item $\displaystyle\lim_{x \to 0}\ \frac{\sqrt{x + 4} - 2}{x}$
        \item $\displaystyle\lim_{x \to 9}\ \frac{\sqrt{x} - 3}{x - 9}$
      \end{enumerate}

% ---------------------------------------------------------
\subsubsection*{One-Sided Limits}

\item Let $f$ be defined by
      \[
        f(x) =
        \begin{cases}
          x^2 + 1 & x < 2 \\
          3       & x = 2 \\
          2x - 1  & x > 2.
        \end{cases}
      \]
      \begin{enumerate}
        \item Find $\displaystyle\lim_{x \to 2^-} f(x)$ and
              $\displaystyle\lim_{x \to 2^+} f(x)$.
        \item Does $\displaystyle\lim_{x \to 2} f(x)$ exist?
              Justify your answer.
        \item Is $f$ continuous at $x = 2$? Why or why not?
      \end{enumerate}

\item For each function, determine whether the two-sided limit
      exists at the indicated point by computing both one-sided
      limits.
      \begin{enumerate}
        \item $f(x) = \dfrac{|x - 1|}{x - 1}$ at $x = 1$
        \item $g(x) = \dfrac{x^2 - 1}{|x - 1|}$ at $x = 1$
      \end{enumerate}

% ---------------------------------------------------------
\subsubsection*{Continuity and Classification of Discontinuities}

\item Determine all points of discontinuity of each function
      and classify each as removable, jump, or infinite.
      \begin{enumerate}
        \item $\displaystyle f(x) = \frac{x^2 - 9}{x - 3}$
        \item $\displaystyle g(x) = \frac{1}{x^2 - 4}$
        \item $\displaystyle h(x) =
              \begin{cases} x + 1 & x \leq 0 \\ x^2 & x > 0 \end{cases}$
      \end{enumerate}

\item For each discontinuity identified below, state which of the
      three conditions for continuity fails (limit existence,
      function defined, limit equals function value).
      \begin{enumerate}
        \item $f(x) = \dfrac{\sin x}{x}$ at $x = 0$
        \item $g(x) = \dfrac{x^2 - 4}{x - 2}$ at $x = 2$
        \item $h(x) = \begin{cases} 1 & x < 0 \\ -1 & x \geq 0 \end{cases}$
              at $x = 0$
        \item $k(x) = \dfrac{1}{(x-1)^2}$ at $x = 1$
      \end{enumerate}

\item Find the value of the constant $k$ that makes $f$ continuous
      everywhere:
      \[
        f(x) =
        \begin{cases}
          kx + 3   & x \leq 2 \\
          x^2 - 1  & x > 2.
        \end{cases}
      \]

\item Determine the interval(s) on which each function is
      continuous.
      \begin{enumerate}
        \item $f(x) = \sqrt{x - 2} + \sqrt{8 - x}$
        \item $g(x) = \dfrac{\tan x}{x^2 + 1}$
        \item $h(x) = \sqrt{\ln(x - 1)}$
      \end{enumerate}

% ---------------------------------------------------------
\subsubsection*{The Squeeze Theorem}

\item Given that $1 - x^2 \leq f(x) \leq \cos(x)$ for all $x$
      near $0$, find $\displaystyle\lim_{x \to 0} f(x)$.

\item Use the Squeeze Theorem to evaluate
      $\displaystyle\lim_{x \to 0}\ x^2 \sin\!\left(\frac{1}{x}\right)$.
      (\textit{Hint: what are the bounds on $\sin(1/x)$?})

% ---------------------------------------------------------
\subsubsection*{The Intermediate Value Theorem}

\item For each function and interval, verify the hypotheses of the
      Intermediate Value Theorem and conclude that a root exists
      in the given interval.
      \begin{enumerate}
        \item $f(x) = x^3 - x - 1$ on $[1, 2]$
        \item $g(x) = e^x - 3x$ on $[0, 1]$
        \item $h(x) = \cos x - x$ on $[0, \pi/2]$
      \end{enumerate}

\item Apply one step of the Bisection Method to
      $f(x) = x^3 - x - 1$ on $[1, 2]$.
      \begin{enumerate}
        \item Evaluate $f$ at the midpoint $d = 1.5$.
        \item Determine which half-interval $[1, 1.5]$ or
              $[1.5, 2]$ is guaranteed to contain a root, and
              state why.
        \item Perform a second bisection step and identify the
              new interval.
      \end{enumerate}

\item Give an example of a function $f$ on $[0, 1]$ such that
      $f(0) < 0$ and $f(1) > 0$, yet $f$ has no root in $(0, 1)$.
      Explain why this does not contradict the IVT.

% ---------------------------------------------------------
\subsubsection*{Limits Involving Infinity}

\item Evaluate the following limits at infinity.
      \begin{enumerate}
        \item $\displaystyle\lim_{x \to \infty}\ \frac{3x^2 - x + 4}{5x^2 + 2}$
        \item $\displaystyle\lim_{x \to \infty}\ \frac{7x^3 + x}{2x^4 - 1}$
        \item $\displaystyle\lim_{x \to -\infty}\ \frac{4x^2 + 1}{x + 3}$
        \item $\displaystyle\lim_{x \to \infty}\ \frac{\sqrt{9x^2 + 1}}{x + 1}$
      \end{enumerate}

\item Find all vertical and horizontal asymptotes of each function.
      \begin{enumerate}
        \item $\displaystyle f(x) = \frac{2x}{x^2 - 1}$
        \item $\displaystyle g(x) = \frac{x^2 - 4}{x^2 + 4}$
        \item $\displaystyle h(x) = \frac{3x^2 - 12}{x - 2}$
      \end{enumerate}

\item Classify each of the following expressions as
      \emph{indeterminate} or \emph{determinate} as $x \to c$
      for the given $c$, and for those that are indeterminate,
      name the form.
      \begin{enumerate}
        \item $\dfrac{x^2 - 1}{x - 1}$ as $x \to 1$
        \item $\dfrac{x + 1}{x - 1}$ as $x \to 1$
        \item $x \ln(x)$ as $x \to 0^+$
        \item $(1 + 1/x)^x$ as $x \to \infty$
      \end{enumerate}

\end{enumerate}

\clearpage
\fi
